% Résumé du mémoire.
% Abstract in French.
%
\chapter*{RÉSUMÉ}\thispagestyle{headings}
\addcontentsline{toc}{compteur}{RÉSUMÉ}

Les réseaux de capteurs sans fil (RCsF) déployés dans des environnements sujets aux catastrophes constituent une infrastructure critique pour la surveillance en temps réel des phénomènes naturels tels que les séismes, les inondations et les glissements de terrain. Cependant, trois défis fondamentaux limitent leur déploiement à grande échelle~: (i)~la durée de vie limitée imposée par les batteries non rechargeables, (ii)~l'empreinte carbone du cycle de vie des dispositifs massivement distribués, et (iii)~la vulnérabilité des politiques de routage apprises par intelligence artificielle face aux attaques et aux hallucinations de modèle. Cette thèse, organisée sous forme de thèse par articles, propose trois cadres complémentaires qui abordent séquentiellement ces défis dans le contexte des réseaux 6G natifs en intelligence artificielle.

Le premier article introduit GAPF (Graph-Attentive Policy Fusion), un cadre méta-\\apprentissage léger (1\,473 paramètres) qui fusionne des politiques expertes gelées d'apprentissage par renforcement multi-agent (QMIX, QTRAN) à travers un encodeur de réseau de neurones graphiques convolutionnels (GCN) à 2 couches et un contrôleur CAS (Contextual Algorithm Selection). L'architecture bi-niveau combine un réseau d'attention monotone Q$\cdot$K pour la scalarisation des récompenses en boucle interne avec des stratégies évolutionnaires pour le réseau d'attention et REINFORCE pour le méta-contrôleur GCN+CAS en boucle externe. Les résultats expérimentaux démontrent un taux de livraison de paquets (PDR) $\geq 98\%$, une consommation d'énergie 2 à 3.5 fois inférieure, et une empreinte mémoire 13 fois moindre que les références, tout en maintenant une inférence en temps quasi-réel.

Le deuxième article présente CALASH (Carbon-Aware Lifecycle-Autonomous Self-Healing), le premier protocole de routage pour RCsF minimisant conjointement l'intensité carbone du cycle de vie (LCI) --- définie comme le CO$_2$ incorporé, opérationnel et de fin de vie par paquet utile livré --- à travers quatre piliers intégrés~: (i)~CADR, qui module le ratio de détection comprimée en fonction de l'intensité carbone du réseau électrique~; (ii)~CARE, un routeur hybride Lyapunov--DQN avec garanties prouvables de conformité au budget carbone~; (iii)~SHDR, une boucle d'auto-guérison autonomique MAPE-K~; et (iv)~LSE, un moteur de durabilité conforme à ISO~14040. Opérant sur une couche physique 6G bi-bande (sub-THz / sub-6\,GHz), CALASH atteint 60\% de durée de vie plus longue que LEACH (1\,486 vs.\ 931 tours), 82.9\% de PDR et 33\% de LCI inférieur (8.99 vs.\ 13.42 gCO$_2$eq/pkt).

Le troisième article propose CASTER-ZT (Zero-Trust Shielded Graph-Structured Decision Control), le premier cadre de contrôle de décision zéro-confiance pour les réseaux de détection de surveillance des catastrophes. CASTER-ZT couple un encodeur GraphSAGE à 2 couches, une politique de récupération apprise, un évaluateur neuronal double-hypothèse de confiance-risque, une quantification d'incertitude par MC-Dropout et un bouclier déterministe produisant des décisions graduées (autoriser, réduire la portée, différer, escalader, bloquer). Dix propositions formelles établissent le conservatisme monotone, l'exclusion d'actuation non autorisée et les garanties de couverture conforme. Une campagne de 2\,420 exécutions démontre 74.2--98.3\% de détection de politiques voyous avec zéro faux positif, un score de récupération pondéré $\omega_{\text{rec}} = 0.733$ (3.7--5.2$\times$ supérieur aux références) et une latence de 3.07\,ms par décision, compatible avec le budget de 10\,ms du near-RT RIC O-RAN.

Collectivement, les trois articles établissent une pile complète --- efficacité énergétique, durabilité carbone et sécurité zéro-confiance --- pour le routage intelligent dans les réseaux de capteurs de surveillance des catastrophes intégrés au 6G.

